% ==============================================================================
%  MSc Thesis – Time Prediction for Embryo Development Using Diffusion Models
%  Author : Ioannis Kasionis
%  Date   : May 2026
% ==============================================================================
\documentclass[12pt, a4paper, twoside]{report}

% ---------------------------------------------------------------------------
%  Packages
% ---------------------------------------------------------------------------
\usepackage[T1]{fontenc}
\usepackage[utf8]{inputenc}
\usepackage[english]{babel}
\usepackage{lmodern}
\usepackage{microtype}

% Page layout
\usepackage[top=3cm, bottom=3cm, left=3.5cm, right=2.5cm,
            headheight=15pt]{geometry}
\usepackage{fancyhdr}
\usepackage{emptypage}   % suppress headers/footers on truly empty pages

% Math
\usepackage{amsmath, amssymb, amsthm, mathtools, bm}

% Graphics / figures
\usepackage{graphicx}
\usepackage{float}
\usepackage{subcaption}
\usepackage[table]{xcolor}
\usepackage{tikz}
\usetikzlibrary{arrows.meta, positioning, shapes.geometric, fit, calc}

% Tables
\usepackage{booktabs}
\usepackage{array}
\usepackage{multirow}
\usepackage{tabularx}

% Captions
\usepackage[font=small, labelfont=bf]{caption}

% Code / algorithms
\usepackage{algorithm}
\usepackage{algpseudocode}
\usepackage{listings}
\lstset{
  basicstyle=\ttfamily\footnotesize,
  breaklines=true,
  frame=single,
  captionpos=b
}

% Hyperlinks (load near last)
\usepackage[hidelinks, colorlinks=true,
            linkcolor=blue!60!black,
            citecolor=green!40!black,
            urlcolor=blue!70!black]{hyperref}
\usepackage{cleveref}

% Bibliography
\usepackage[backend=biber, style=ieee, sortcites=true]{biblatex}
\addbibresource{references.bib}

% Misc
\usepackage{setspace}
\usepackage{parskip}
\usepackage{enumitem}
\usepackage{tocloft}
\usepackage{acronym}
\usepackage{rotating}
\usepackage{pdfpages}
\usepackage{lipsum}   % remove in final version

% ---------------------------------------------------------------------------
%  Custom commands
% ---------------------------------------------------------------------------
\newcommand{\thesis}{Time Prediction for Embryo Development Using Diffusion Models}
\newcommand{\theauthor}{Ioannis Kasionis}
\newcommand{\theyear}{2026}

\newcommand{\x}{\mathbf{x}}
\newcommand{\eps}{\boldsymbol{\varepsilon}}
\newcommand{\calL}{\mathcal{L}}
\newcommand{\calN}{\mathcal{N}}
\newcommand{\ctx}{\text{ctx}}
\newcommand{\R}{\mathbb{R}}
\newcommand{\E}{\mathbb{E}}

\DeclareMathOperator{\SNR}{SNR}

% ---------------------------------------------------------------------------
%  Page style
% ---------------------------------------------------------------------------
\pagestyle{fancy}
\fancyhf{}
\fancyhead[LE]{\nouppercase{\leftmark}}
\fancyhead[RO]{\nouppercase{\rightmark}}
\fancyfoot[C]{\thepage}
\renewcommand{\headrulewidth}{0.4pt}

\setstretch{1.3}

% ==============================================================================
\begin{document}

% ---------------------------------------------------------------------------
%  Front matter
% ---------------------------------------------------------------------------
\pagenumbering{roman}

% ---- Title page ----
\begin{titlepage}
\begin{center}
  \vspace*{1cm}
  {\LARGE \textbf{Department of Computer Science and Engineering}}\\[0.5em]
  {\large MSc in Artificial Intelligence}\\
  \vspace{2cm}
  \includegraphics[width=0.25\textwidth]{figures/university_logo}\\  % replace with your logo
  \vspace{2cm}
  {\Huge \textbf{\thesis}}\\
  \vspace{1.5cm}
  {\large Master's Thesis}\\
  \vspace{1cm}
  {\Large \theauthor}\\
  \vspace{2cm}
  \begin{tabular}{ll}
    \textbf{Supervisor:} & [Supervisor Name] \\
  \end{tabular}
  \vfill
  {\large \theyear}
\end{center}
\end{titlepage}

\cleardoublepage

% ---- Declaration ----
\chapter*{Declaration}
\addcontentsline{toc}{chapter}{Declaration}
I hereby declare that this thesis is my own work and that, to the best of my
knowledge and belief, it contains no material previously published or written
by another person, nor material which has been accepted for the award of any
other degree or diploma, except where due acknowledgement has been made.

\bigskip
\noindent Ioannis Kasionis \hfill \theyear

\cleardoublepage

% ---- Abstract ----
\chapter*{Abstract}
\addcontentsline{toc}{chapter}{Abstract}

This thesis investigates \emph{temporal next-frame prediction} of human embryo
development using conditional diffusion models.
Given a short temporal context window of $k$ consecutive embryo frames ending
at time $t$, the objective is to synthesize the subsequent frame at $t{+}1$
and, by iteratively rolling the context window forward, generate arbitrarily
long future sequences.

The work covers the full pipeline from raw data to trained model:
an automated preprocessing stage applies embryo presence filtering, embryo-centred
circular cropping, and contrast-limited adaptive histogram equalisation (CLAHE);
a binary presence classifier removes invalid frames before training;
a multi-class morpho-kinetic phase classifier ($16$ developmental phases) is
trained for analysis and optional classifier-guided rollout;
and a context-conditioned denoising diffusion probabilistic model (DDPM) is
trained to predict the next frame via a UNet equipped with cross-attention over
context tokens.

Quantitative evaluation on a held-out embryo-level test split reports an
average SSIM of $0.436$ and LPIPS of $0.116$ for one-step prediction.
Multi-step free-running rollouts demonstrate plausible short-horizon generation
while exhibiting gradual structural drift over longer horizons, consistent with
autoregressive error accumulation.
The results establish a strong diffusion-based baseline for embryo time-lapse
prediction and motivate future work on morphologically consistent long-rollout
generation.

\bigskip
\noindent\textbf{Keywords:} diffusion models, embryo development, next-frame
prediction, morpho-kinetics, time-lapse microscopy, conditional generation.

\cleardoublepage

% ---- Acknowledgements ----
\chapter*{Acknowledgements}
\addcontentsline{toc}{chapter}{Acknowledgements}

I would like to thank my supervisor and all those who supported this work.
The embryo time-lapse dataset was provided by Gomez et al.~\cite{gomez2022}
under the CC BY 4.0 licence.

\cleardoublepage

% ---- Table of Contents / Lists ----
\tableofcontents
\cleardoublepage

\listoffigures
\addcontentsline{toc}{chapter}{List of Figures}
\cleardoublepage

\listoftables
\addcontentsline{toc}{chapter}{List of Tables}
\cleardoublepage

% ---- Abbreviations ----
\chapter*{List of Abbreviations}
\addcontentsline{toc}{chapter}{List of Abbreviations}
\begin{acronym}[CLAHE]
  \acro{DDPM}{Denoising Diffusion Probabilistic Model}
  \acro{DDIM}{Denoising Diffusion Implicit Model}
  \acro{CNN}{Convolutional Neural Network}
  \acro{UNet}{U-shaped Network}
  \acro{CLAHE}{Contrast Limited Adaptive Histogram Equalisation}
  \acro{MSE}{Mean Squared Error}
  \acro{MAE}{Mean Absolute Error}
  \acro{SSIM}{Structural Similarity Index Measure}
  \acro{LPIPS}{Learned Perceptual Image Patch Similarity}
  \acro{SNR}{Signal-to-Noise Ratio}
  \acro{EMA}{Exponential Moving Average}
  \acro{AMP}{Automatic Mixed Precision}
  \acro{AdamW}{Adam with Weight Decay}
  \acro{CFG}{Classifier-Free Guidance}
  \acro{IVF}{In Vitro Fertilisation}
\end{acronym}

\cleardoublepage

% ---------------------------------------------------------------------------
%  Main matter
% ---------------------------------------------------------------------------
\pagenumbering{arabic}

% ==============================================================================
\chapter{Introduction}
\label{chap:introduction}
% ==============================================================================

\section{Motivation}
\label{sec:motivation}

Human embryo development is a complex, time-dependent biological process that
exhibits characteristic morphological transitions across well-defined
developmental phases.
In the context of assisted reproduction, the selection of the most viable
embryo for transfer is a critical clinical decision that directly influences
the probability of successful implantation.
Modern time-lapse incubators continuously acquire microscopy images of
developing embryos at fixed temporal intervals, creating high-resolution
recordings of the entire pre-implantation period.
These recordings encode rich temporal and morphological information that is
currently interpreted by trained embryologists who score embryos according to
standardised morpho-kinetic criteria.

Machine learning methods have begun to complement manual assessment by
automating classification of developmental phases and grading of embryo
quality.
However, most existing approaches treat the problem as a static image
classification task, ignoring the inherent temporal dynamics of embryo
development.
Generative models that can plausibly \emph{predict} future developmental
states could open new clinical and scientific possibilities: prospective
viability assessment, simulation-based embryo selection, and data augmentation
for downstream classification models.

\section{Problem Statement}
\label{sec:problem}

This thesis addresses the problem of \textbf{next-frame prediction} in human
embryo time-lapse sequences.
Formally, given a context window of $k$ consecutive, preprocessed grayscale
embryo frames
\[
  \mathbf{c}_t = \bigl(x_{t-k+1}, \, x_{t-k+2}, \, \ldots, \, x_t\bigr)
\]
observed up to time $t$, the goal is to synthesise a plausible frame
$\hat{x}_{t+1}$ at the next time point.
Multi-step rollouts are obtained by iteratively appending the generated frame
to the context and advancing the window:
\[
  \hat{x}_{t+j+1} = f\!\left(\hat{x}_{t+j-k+1}, \ldots, \hat{x}_{t+j}\right),
  \quad j = 0, 1, 2, \ldots
\]

Three desiderata guide the approach:
\begin{enumerate}[label=(\roman*)]
  \item \textbf{Visual realism}: generated frames must resemble authentic
        microscopy images in texture and intensity distribution.
  \item \textbf{Temporal consistency}: generated frames must be plausible
        continuations of the conditioning context.
  \item \textbf{Morphological correctness}: generated frames should preserve
        and appropriately evolve coarse embryo structure across phase
        transitions.
\end{enumerate}

\section{Approach Overview}
\label{sec:approach}

The approach follows a structured pipeline comprising four components:
\begin{enumerate}
  \item \textbf{Preprocessing}: an automated pipeline converts raw bronze
        frames into a filtered, normalised \emph{silver} dataset.
  \item \textbf{Presence classifier}: a binary convolutional model filters
        frames that do not contain a visible embryo before training.
  \item \textbf{Phase classifier}: a multi-class convolutional model maps
        single frames to one of 16 morpho-kinetic phases, used for analysis
        and optional guided rollout.
  \item \textbf{Diffusion model}: a context-conditioned
        \texttt{UNet2DConditionModel} is trained for noise prediction on
        strictly adjacent frame pairs.
\end{enumerate}

\section{Contributions}
\label{sec:contributions}

The main contributions of this thesis are:
\begin{itemize}
  \item A fully automated, reproducible preprocessing pipeline for the public
        Gomez et al.\ embryo time-lapse dataset that produces a high-quality
        silver dataset across seven focal planes.
  \item A morpho-kinetic phase classifier achieving strong performance across
        16 developmental phases.
  \item A context-conditioned diffusion model with absolute frame-index
        temporal embedding and cross-attention context conditioning, designed
        specifically for embryo next-frame prediction.
  \item A comprehensive quantitative and qualitative evaluation covering
        one-step prediction, transition vs.\ non-transition analysis, and
        multi-step free-running rollouts.
\end{itemize}

\section{Thesis Structure}
\label{sec:structure}

The remainder of this thesis is organised as follows.
\Cref{chap:background} reviews relevant background on diffusion models and
related generative approaches.
\Cref{chap:dataset} describes the dataset and the preprocessing pipeline.
\Cref{chap:auxiliary} presents the auxiliary classification models.
\Cref{chap:diffusion} introduces the diffusion model architecture, training
objective, and sampling procedure.
\Cref{chap:experiments} presents experimental results.
\Cref{chap:discussion} discusses limitations and future directions.
\Cref{chap:conclusion} concludes the thesis.

% ==============================================================================
\chapter{Background and Related Work}
\label{chap:background}
% ==============================================================================

\section{Diffusion Probabilistic Models}
\label{sec:diffusion_background}

Diffusion probabilistic models~\cite{sohl2015deep,ho2020ddpm} are a class of
latent variable models that define a \emph{forward process} that gradually
corrupts data $x_0$ into Gaussian noise over $T$ discrete timesteps, and a
\emph{reverse process} that learns to invert this corruption.

\subsection{Forward Process}
\label{ssec:forward}

The forward process is a fixed Markov chain:
\[
  q(x_{1:T} \mid x_0) = \prod_{t=1}^{T} q(x_t \mid x_{t-1}), \qquad
  q(x_t \mid x_{t-1}) = \calN\!\left(x_t;\, \sqrt{1-\beta_t}\, x_{t-1},\, \beta_t \mathbf{I}\right),
\]
where $\beta_1 < \beta_2 < \cdots < \beta_T$ is a variance schedule.
Defining $\alpha_t = 1 - \beta_t$ and $\bar{\alpha}_t = \prod_{s=1}^{t} \alpha_s$,
one can sample $x_t$ directly from $x_0$ in closed form:
\[
  x_t = \sqrt{\bar{\alpha}_t}\, x_0 + \sqrt{1-\bar{\alpha}_t}\, \eps,
  \qquad \eps \sim \calN(\mathbf{0}, \mathbf{I}).
\]

\subsection{Reverse Process and Training Objective}
\label{ssec:reverse}

The reverse process is parameterised by a neural network
$\eps_\theta(x_t, t)$ that predicts the noise component $\eps$:
\[
  \calL_{\text{simple}} = \E_{x_0,\, t,\, \eps}\!\left[
    \left\|\eps - \eps_\theta(x_t, t)\right\|^2
  \right].
\]
Ho et al.~\cite{ho2020ddpm} showed that this simplified objective
outperforms the full variational lower bound in practice.

\subsection{Min-SNR Reweighting}
\label{ssec:min_snr}

The signal-to-noise ratio at timestep $t$ is
$\text{SNR}(t) = \bar\alpha_t / (1 - \bar\alpha_t)$.
Hang et al.~\cite{hang2023minsnr} propose to clip the per-timestep loss
weight at $\gamma_{\text{SNR}} = 5$:
\[
  w(t) = \min\!\left(\text{SNR}(t),\; \gamma_{\text{SNR}}\right),
\]
which prevents the training loss from being dominated by high-noise
timesteps where learning the signal structure is difficult.

\subsection{DDIM Sampling}
\label{ssec:ddim}

Denoising Diffusion Implicit Models~\cite{song2021ddim} define a
non-Markovian inference process that produces the same marginals as the DDPM
forward process but allows sampling in far fewer steps.
The deterministic ($\eta = 0$) DDIM update is:
\begin{align}
  \hat{x}_0^{(t)} &= \frac{x_t - \sqrt{1-\bar\alpha_t}\,\eps_\theta(x_t,t)}
                        {\sqrt{\bar\alpha_t}}, \\
  x_{t-1} &= \sqrt{\bar\alpha_{t-1}}\, \hat{x}_0^{(t)}
             + \sqrt{1-\bar\alpha_{t-1}}\, \eps_\theta(x_t, t).
\end{align}

\subsection{Conditional Diffusion Models}
\label{ssec:conditional}

Conditioning information $\mathbf{c}$ is incorporated by replacing
$\eps_\theta(x_t, t)$ with $\eps_\theta(x_t, t, \mathbf{c})$.
Common conditioning mechanisms include class embedding concatenation,
cross-attention over token sequences~\cite{rombach2022ldm,ramesh2022dalle2},
and channel-wise concatenation of conditioning images.
In this work the context frames are encoded into a spatial token sequence
that is consumed by a cross-attention module inside the UNet denoiser.

\section{Convolutional Architectures for Embryo Analysis}
\label{sec:related_cnn}

Several works have applied CNNs to embryo quality assessment and phase
classification.
Khosravi et al.~\cite{khosravi2019} demonstrate that deep learning models
trained on time-lapse data can predict morpho-kinetic parameters with
accuracy comparable to experienced embryologists.
Tran et al.~\cite{tran2019} train recurrent convolutional models to
predict implantation outcome from morpho-kinetic annotations.
Our phase classifier follows the same single-frame classification paradigm
but extends it to 16 fine-grained developmental phases.

\section{Generative Models for Medical Image Synthesis}
\label{sec:related_generative}

Generative adversarial networks (GANs) have been widely used for medical image
synthesis, including augmentation of training sets and simulating imaging
artefacts~\cite{yi2019review}.
More recently, diffusion models have emerged as the leading paradigm for
high-fidelity image generation, with stable training and improved coverage
compared to GANs~\cite{dhariwal2021beatgans,ho2020ddpm}.
Their application to time-series and video prediction is an active area of
research~\cite{ho2022video,harvey2022flexible}, and this work adapts these
ideas to the embryo domain.

% ==============================================================================
\chapter{Dataset and Preprocessing}
\label{chap:dataset}
% ==============================================================================

\section{Dataset}
\label{sec:dataset}

The experiments use the public time-lapse embryo dataset introduced by
Gomez et al.~\cite{gomez2022}.
The dataset contains recordings of \textbf{704 human embryos} from IVF
procedures.
For each embryo, a sequence of JPEG images ($500 \times 500$ pixels) is
provided together with a phase annotation CSV that records, for each
developmental phase, the first frame at which the transition was observed.

The raw dataset is organised as follows:
\begin{itemize}
  \item \texttt{data/embryo\_dataset\_bronze/} — raw per-embryo JPEG frames,
        organised by focal plane.
  \item \texttt{data/embryo\_dataset\_annotations/} — one
        \texttt{<embryo\_id>\_phases.csv} per embryo.
\end{itemize}

Seven focal planes are available (F$-45$, F$-30$, F$-15$, F$0$, F$15$, F$30$,
F$45$), where F$0$ refers to the focus plane closest to the equatorial section
of the embryo.
All primary experiments use the F$0$ plane unless stated otherwise.

\subsection{Developmental Phases}
\label{ssec:phases}

Each embryo is annotated with up to 16 morpho-kinetic phases
(Table~\ref{tab:phases}), spanning from the second polar body extrusion
(\texttt{pPB2}) to the full blastocyst stage (\texttt{pHB}).
Phase transitions are irreversible and strictly ordered in time.

\begin{table}[H]
\centering
\caption{The 16 morpho-kinetic phase labels used in this work.}
\label{tab:phases}
\begin{tabular}{lll}
\toprule
\textbf{ID} & \textbf{Label} & \textbf{Description} \\
\midrule
1  & pPB2  & Second polar body extrusion \\
2  & pPNa  & Pronuclear appearance \\
3  & pPNf  & Pronuclear fading \\
4  & p2    & 2-cell stage \\
5  & p3    & 3-cell stage \\
6  & p4    & 4-cell stage \\
7  & p5    & 5-cell stage \\
8  & p6    & 6-cell stage \\
9  & p7    & 7-cell stage \\
10 & p8    & 8-cell stage \\
11 & p9+   & ${\geq}9$-cell stage \\
12 & pM    & Morula \\
13 & pSB   & Start of blastulation \\
14 & pB    & Blastocyst \\
15 & pEB   & Expanded blastocyst \\
16 & pHB   & Hatching/hatched blastocyst \\
\bottomrule
\end{tabular}
\end{table}

\section{Preprocessing Pipeline}
\label{sec:preprocessing}

\subsection{Design Principles}
\label{ssec:preprocess_principles}

The preprocessing pipeline converts raw bronze frames into a high-quality
\emph{silver} dataset.
The design adheres to three principles:
\begin{enumerate}
  \item \textbf{Determinism}: the pipeline is fully automated and reproducible.
        Running it twice on the same input produces bit-identical output.
  \item \textbf{Index preservation}: output frames retain their original
        1-based frame indices (e.g., \texttt{00001.jpeg}, \texttt{00037.jpeg}).
        Frames removed by filtering create explicit index gaps; there is no
        renumbering.
  \item \textbf{No temporal smoothing}: no interpolation or temporal
        averaging is applied.
\end{enumerate}

\subsection{Embryo Presence Filtering}
\label{ssec:presence_filtering}

A binary presence classifier (described in \Cref{sec:presence_classifier}) is
applied to each frame independently.
Frames for which the classifier predicts \emph{no visible embryo} are discarded.
This step removes background-only frames, acquisition artefacts, and severely
degraded images before any further processing occurs.

Frames are first sub-sampled according to a fixed temporal stride before
classifier inference to reduce computational cost.

\subsection{Spatial Cropping}
\label{ssec:cropping}

A circular crop centred on the embryo is applied using Hough circle detection.
The Hough transform is run on a Gaussian-blurred edge map to robustly localise
the circular well boundary that surrounds each embryo.
If circle detection fails (e.g.\ due to overexposed frames or detection
ambiguity), a fall-back centre crop of the same target size is used.
The resulting crop is then resized to a fixed resolution.

\subsection{Contrast Normalisation}
\label{ssec:clahe}

Contrast Limited Adaptive Histogram Equalisation (CLAHE) is applied after
resizing.
CLAHE subdivides the image into a grid of tiles and equalises the histogram
within each tile while limiting contrast amplification via a configurable clip
limit.
This reduces illumination variability across embryos and focal planes, and
enhances morphological detail that may otherwise be obscured by uneven
background illumination.

\subsection{Output Organisation}
\label{ssec:silver_layout}

Processed frames are written to:
\begin{verbatim}
data/embryo_dataset_silver/<plane_id>/<embryo_id>/<frame_idx>.jpeg
\end{verbatim}
Phase annotation CSV files are updated to reflect the retained frame indices
after filtering and are copied alongside the images.

\subsection{Quality Control}
\label{ssec:quality_control}

Basic consistency checks are applied throughout the pipeline:
file readability is verified before processing, and frames that fail to load
or produce invalid outputs are skipped with a warning.
Misclassifications by the presence model may introduce short temporal gaps.
These gaps are preserved and handled explicitly during training pair
construction, rather than corrected heuristically.

\section{Data Summary}
\label{sec:data_summary}

After preprocessing, the silver dataset contains:
\begin{itemize}
  \item \textbf{704 embryos} across seven focal planes (F$-45$ to F$45$).
  \item Variable-length sequences per embryo, reflecting natural biological
        variation in developmental rate.
  \item Index gaps where frames were filtered by the presence classifier.
\end{itemize}
All experiments in this thesis use the F$0$ focal plane.

% ==============================================================================
\chapter{Auxiliary Classification Models}
\label{chap:auxiliary}
% ==============================================================================

Two auxiliary convolutional classification models support the pipeline:
an embryo presence classifier and a morpho-kinetic phase classifier.
Neither model participates in the diffusion training objective.

\section{Embryo Presence Classifier}
\label{sec:presence_classifier}

\subsection{Task and Architecture}
\label{ssec:presence_arch}

The presence classifier is a binary convolutional neural network that receives
a single grayscale embryo frame and predicts whether it contains a visible,
well-formed embryo (\texttt{T=1}) or not (\texttt{T=0}).

Frames may be labelled negative for several reasons: the embryo well is empty,
the acquisition was corrupted, the frame was captured before or after the
embryo was present in the field of view, or imaging conditions were
sufficiently degraded that morphological assessment is impossible.

\subsection{Training}
\label{ssec:presence_training}

The model is trained on manually curated examples of valid and invalid frames
drawn from the bronze dataset.
Binary cross-entropy loss is used.
Standard data augmentations (horizontal and vertical flips, brightness and
contrast jitter) are applied during training.

\subsection{Role in the Pipeline}
\label{ssec:presence_role}

The presence classifier is applied during preprocessing and determines which
frames are retained in the silver dataset.
It is not used during diffusion model training or sampling.

\section{Morpho-Kinetic Phase Classifier}
\label{sec:phase_classifier}

\subsection{Task}
\label{ssec:phase_task}

The phase classifier assigns a single embryo frame to one of the 16
developmental phases listed in Table~\ref{tab:phases}.
It operates at the frame level and does not exploit temporal context.

\subsection{Architecture and Training}
\label{ssec:phase_arch}

The classifier follows a convolutional architecture with grayscale input and
a 16-way softmax output.
An EfficientNet-based backbone is used, pre-trained on ImageNet and
fine-tuned on the embryo data.
Timestep-dependent Gaussian noise is added to input frames during training to
improve robustness to mild imaging variability.

Ground-truth labels are derived from the phase annotation CSV files by
assigning each frame the phase in whose annotated interval it falls.
Significant class imbalance exists (e.g.\ early 2-cell and blastocyst phases
are under-represented), which is mitigated using weighted random sampling.

Standard augmentations — horizontal flips, brightness and contrast jitter,
and small rotations — are applied during training, while validation and
testing use deterministic preprocessing.

\subsection{Performance}
\label{ssec:phase_results}

Table~\ref{tab:phase_results} summarises the classifier's performance on the
held-out test set.

\begin{table}[H]
\centering
\caption{Phase classifier performance on the test set.}
\label{tab:phase_results}
\begin{tabular}{lc}
\toprule
\textbf{Metric} & \textbf{Value} \\
\midrule
Weighted F1 (F0 plane) & 0.643 \\
\bottomrule
\end{tabular}
\end{table}

The confusion matrix (Figure~\ref{fig:confusion_matrix}) reveals that
adjacent phases are most commonly confused, which is expected given that
phase transitions occur gradually and the boundary between consecutive phases
is not always visually unambiguous in a single frame.

\begin{figure}[H]
\centering
% \includegraphics[width=0.75\textwidth]{figures/confusion_matrix}
\caption{Row-normalised confusion matrix of the phase classifier on the test
         set. Adjacent-phase confusion dominates.}
\label{fig:confusion_matrix}
\end{figure}

\subsection{Role in the Diffusion Pipeline}
\label{ssec:phase_role}

The phase classifier is \emph{not} used during diffusion model training.
At sampling time it can optionally steer multi-step rollouts via a monotonic
adjacent-confidence phase policy, which advances the phase label only when
the classifier confidence in the next adjacent phase exceeds a threshold
($\tau = 0.55$ in the default configuration).
This guidance is disabled in all quantitative evaluations reported in this
thesis.

\begin{table}[H]
\centering
\caption{Summary of auxiliary model usage in the full pipeline.}
\label{tab:auxiliary_summary}
\begin{tabular}{lcc}
\toprule
\textbf{Model} & \textbf{Used in training} & \textbf{Used in sampling} \\
\midrule
Presence classifier & Yes (preprocessing) & No \\
Phase classifier    & No                  & No (analysis only) \\
\bottomrule
\end{tabular}
\end{table}

% ==============================================================================
\chapter{Diffusion Model for Next-Frame Prediction}
\label{chap:diffusion}
% ==============================================================================

\section{Task Formulation}
\label{sec:diff_task}

The diffusion model is trained for \textbf{next-frame prediction} on the F$0$
silver dataset.
Let $x_{t+1} \in \R^{H \times W}$ denote the target frame and
$\mathbf{c}_t = (x_{t-k+1}, \ldots, x_t)$ the context window of the $k$
most recent frames.
The model learns the conditional distribution:
\[
  p_\theta(x_{t+1} \mid \mathbf{c}_t, \, n_{t+1}),
\]
where $n_{t+1}$ is the absolute 1-based frame index of the target.

\section{Data Preparation}
\label{sec:diff_data}

\subsection{Frame Index and Phase Assignment}
\label{ssec:diff_index}

A cached frame index is built by scanning the silver dataset.
Frame numbers are parsed from filenames (e.g.\ \texttt{00042.jpeg}
$\to$ frame index $42$) and treated as absolute 1-based time points.
Each frame is assigned a phase label derived from the corresponding
\texttt{<embryo\_id>\_phases.csv} by finding the phase interval that contains
the frame index.
Only frames with a valid phase assignment are included in the index.

\subsection{Dataset Splits}
\label{ssec:splits}

Train/validation/test splits are performed at the \textbf{embryo level} to
prevent data leakage: all frames from a given embryo appear in exactly one
split.
Splits are stratified by the dominant phase label per embryo to maintain
similar phase distributions.
The split sizes are given in Table~\ref{tab:splits}.

\begin{table}[H]
\centering
\caption{Embryo-level split sizes (cached to disk).}
\label{tab:splits}
\begin{tabular}{lrr}
\toprule
\textbf{Split} & \textbf{Embryos} & \textbf{Fraction} \\
\midrule
Train      & 565 & 80\% \\
Validation & 69  & 10\% \\
Test       & 69  & 10\% \\
\bottomrule
\end{tabular}
\end{table}

\subsection{Pair Construction}
\label{ssec:pair_construction}

Training samples are constructed as deterministic, overlapping windows.
A valid window $(\mathbf{c}_t, x_{t+1})$ requires that all $k+1$ consecutive
frames exist in the index.
Missing frames (filtered by the presence classifier during preprocessing)
cause windows to be skipped rather than reindexed:
\[
  \text{valid pair} \iff
    \{t{-}k{+}1, t{-}k{+}2, \ldots, t, t{+}1\} \subseteq \text{retained frames}.
\]

For validation and test sets, each pair is additionally labelled as a
\emph{transition} if the phase label changes between $t$ and $t+1$, or
\emph{non-transition} otherwise.
This labelling enables separate qualitative and quantitative analysis of
cross-phase and within-phase prediction.

\section{Model Architecture}
\label{sec:architecture}

The model consists of two components: a \textbf{context encoder} and a
\textbf{UNet denoiser}.
Figure~\ref{fig:architecture} gives a schematic overview.

\begin{figure}[H]
\centering
% \includegraphics[width=0.85\textwidth]{figures/architecture}
\caption{Architecture overview. The context encoder produces a sequence of
         spatial tokens that condition the UNet denoiser via cross-attention.
         The noisy target, a compressed context summary, and an optional
         absolute frame-index map are concatenated as UNet input channels.}
\label{fig:architecture}
\end{figure}

\subsection{Context Encoder}
\label{ssec:ctx_encoder}

Each of the $k$ context frames is processed independently by a lightweight
shared CNN:
\begin{align*}
  &\text{Conv}(1 \to 32, 3{\times}3, s2) \to \text{GELU} \\
  &\text{Conv}(32 \to 64, 3{\times}3, s2) \to \text{GELU} \\
  &\text{Conv}(64 \to D, 3{\times}3, s1) \to \text{GELU},
\end{align*}
where $D = 96$ is the context embedding dimension.
The output feature map is resized to a fixed spatial grid of
$16 \times 16$ tokens via bilinear interpolation.

Three additive positional embeddings are added to each spatial feature map:
\begin{enumerate}
  \item A \emph{context position embedding} $\mathbf{p}_j \in \R^D$ that
        encodes the temporal order of the $j$-th context frame ($j=1,\ldots,k$).
  \item A \emph{spatial position embedding}
        $\mathbf{P}_{\text{spatial}} \in \R^{D \times 16 \times 16}$
        (a learnable parameter shared across all frames and timesteps).
  \item An \emph{absolute frame-index embedding}
        $\text{Emb}(n_j) \in \R^D$ derived from a learned embedding table,
        where $n_j$ is the 1-based frame index of the $j$-th context frame.
\end{enumerate}

After adding all positional embeddings, the $k$ feature maps are flattened and
concatenated into a token sequence of length $k \times 16 \times 16$,
which is passed to the UNet via cross-attention:
\[
  \mathbf{Z} \in \R^{k \cdot 16 \cdot 16 \times D}.
\]

\subsection{Context Summary Channel}
\label{ssec:ctx_summary}

In addition to the cross-attention token sequence, a compact \emph{context
summary image} is constructed by projecting the $k$ context channels through
a shallow CNN:
\[
  f_\ctx: \R^{k \times H \times W} \to \R^{c_\ctx \times H \times W},
  \qquad c_\ctx = 4.
\]
This summary is concatenated with the noisy target frame as additional input
channels to the UNet, giving the denoiser direct pixel-level access to
recent context information without routing everything through cross-attention.

\subsection{Absolute Frame-Index Time Map}
\label{ssec:time_map}

When \texttt{USE\_TIME\_MAP=True}, the absolute index $n_{t+1}$ of the target
frame is embedded via a learned embedding table of dimension
$d_\text{time} = 8$, projected to a scalar via a linear layer, and broadcast
to a spatial map of size $1 \times H \times W$:
\[
  \mathbf{m}_{t+1} = \text{proj}(\text{Emb}(n_{t+1})) \cdot \mathbf{1}_{H \times W}
  \;\in\; \R^{1 \times H \times W}.
\]
This map is concatenated as a third input channel group, giving the denoiser
explicit information about elapsed developmental time.

\subsection{UNet Denoiser}
\label{ssec:unet}

The denoising backbone is a \texttt{UNet2DConditionModel} from the
\texttt{diffusers} library~\cite{von2022diffusers}, configured with:
\begin{itemize}
  \item $1{+}4{+}1 = 6$ input channels (noisy target + context summary + time map).
  \item 1 output channel (predicted noise on the target frame).
  \item 5 encoder/decoder resolution levels with channel widths
        $(32, 64, 128, 256, 256)$.
  \item Cross-attention at the bottleneck resolution only
        (\texttt{CrossAttnDownBlock2D}, \texttt{CrossAttnUpBlock2D}).
  \item Cross-attention dimension $D = 96$.
  \item Gradient checkpointing enabled to reduce GPU memory.
\end{itemize}

\section{Diffusion Schedule and Prediction}
\label{sec:schedules}

\subsection{Noise Schedule}
\label{ssec:noise_schedule}

A \textbf{linear beta schedule} is used for training:
$\beta_1 = 10^{-4}$, $\beta_T = 2 \times 10^{-2}$, with $T = 1000$ total
timesteps.
This schedule is implemented via the \texttt{DDPMScheduler} from
\texttt{diffusers}.

\subsection{Residual Prediction}
\label{ssec:residual}

Rather than denoising the absolute target frame $x_{t+1}$, the model is
trained to denoise the normalised inter-frame residual:
\[
  r_{t+1} = \frac{x_{t+1} - x_t}{2} \;\in\; [-1, 1],
\]
where $x_t$ is the most recent context frame.
This \emph{residual parameterisation} concentrates the generative task on the
small incremental change between consecutive frames, which is far easier to
learn than the full absolute intensity, and also closes the gap between the
model and the copy-last-frame baseline.
The predicted residual is combined with $x_t$ at sampling time to recover
the full predicted next frame $\hat{x}_{t+1} = x_t + 2\,\hat{r}_{t+1}$.

\subsection{DDIM Sampling}
\label{ssec:ddim_sampling}

Inference uses \textbf{DDIM} with $S = 150$ denoising steps and
$\eta = 0$ (fully deterministic).
This reduces sampling time by a factor of $T/S = 6.7\times$ compared to
full DDPM ancestral sampling.

\section{Training Objective}
\label{sec:training_objective}

The primary training loss is mean squared error between the model's noise
prediction and the true noise:
\begin{equation}
  \calL_{\text{MSE}} = \E_{x_0,\, t,\, \eps}\!\left[
    \left\|\eps - \eps_\theta(x_t, t, \mathbf{Z}, \mathbf{m})\right\|^2
  \right].
  \label{eq:mse_loss}
\end{equation}

\subsection{L1 Regularisation on the Predicted Clean Image}
\label{ssec:l1_loss}

An additional $\ell_1$ penalty is applied on the reconstructed clean image
$\hat{x}_0 = (x_t - \sqrt{1-\bar\alpha_t}\,\eps_\theta) / \sqrt{\bar\alpha_t}$:
\begin{equation}
  \calL_{x_0} = \left\|\hat{x}_0 - x_0\right\|_1.
  \label{eq:l1_loss}
\end{equation}
This term penalises large pixel-level errors on the clean prediction and
provides a complementary gradient signal to the noise-space loss.
The weight is set to $\lambda_{x_0} = 0.02$.

\subsection{Laplacian Edge Loss}
\label{ssec:laplacian_loss}

A Laplacian penalty encourages the model to preserve high-frequency image
detail (edges and texture) that $\ell_2$-on-noise tends to smooth out:
\begin{equation}
  \calL_{\text{Lap}} =
    \left\|\nabla^2 \hat{x}_0 - \nabla^2 x_0\right\|^2_2,
  \label{eq:lap_loss}
\end{equation}
where $\nabla^2$ denotes a fixed $3 \times 3$ discrete Laplacian convolution.
The weight is $\lambda_{\text{Lap}} = 0.05$.

\subsection{Transition-Aware Weighting}
\label{ssec:transition_loss}

To ensure that the model receives adequate supervision at morpho-kinetic phase
transitions, a weighted random sampler is used during training to set the
fraction of transition pairs to $20\%$ of each batch.
An additive transition loss term is applied with weight
$\lambda_{\text{trans}} = 0.5$.

\subsection{Min-SNR Reweighting}
\label{ssec:minsnr}

Min-SNR reweighting~\cite{hang2023minsnr} is applied with
$\gamma_{\text{SNR}} = 5.0$ to balance gradient contributions across the noise
schedule.
The total loss is:
\begin{equation}
  \calL = w(t)\cdot\calL_{\text{MSE}}
        + \lambda_{x_0}\,\calL_{x_0}
        + \lambda_{\text{Lap}}\,\calL_{\text{Lap}},
  \label{eq:total_loss}
\end{equation}
where $w(t) = \min(\SNR(t), \gamma_{\text{SNR}})$.

\section{Context Perturbation}
\label{sec:ctx_perturbation}

A known challenge in autoregressive generation is \emph{exposure bias}: the
model is trained on ground-truth context but evaluated on its own
(potentially erroneous) outputs.
To bridge this gap, context frames are randomly perturbed during training.
A recency-weighted scheme applies perturbation most strongly to the most
recent context frame:
\begin{itemize}
  \item The newest frame ($x_t$) is perturbed with probability
        $p_{\text{perturb}} = 0.50$.
  \item The $j$-th-newest frame is perturbed with probability
        $p \cdot d^{j-1}$ where $d = 0.50$ is the decay factor.
  \item Perturbation consists of additive Gaussian noise
        ($\sigma = 0.07$) followed by brightness scaling
        ($\text{scale} \sim \mathcal{U}[0.80, 1.00]$).
\end{itemize}
The rationale is that rollout errors accumulate in the most recent frames,
so training with perturbed recent frames better represents the distribution
encountered during free-running rollout.

\section{Optimiser and Regularisation}
\label{sec:optimiser}

\subsection{Optimiser}
\label{ssec:optimizer}

The model is trained with \textbf{AdamW}~\cite{loshchilov2019decoupled}
(learning rate $\eta = 10^{-4}$, default $\beta_1, \beta_2$, weight decay
$10^{-2}$).
Gradient norms are clipped to $1.0$ at each step.

\subsection{Learning Rate Schedule}
\label{ssec:lr_schedule}

A \textbf{linear warm-up} over 10 epochs (starting from $0.10 \times \eta$)
is followed by \textbf{cosine annealing} with minimum learning rate
$\eta_{\min} = 10^{-7}$.

\subsection{Exponential Moving Average}
\label{ssec:ema}

An \textbf{EMA} of model weights is maintained with decay $0.9995$,
starting from step 100.
EMA weights are used for all validation, sampling, and checkpoint selection.

\subsection{Data Augmentation}
\label{ssec:augmentation}

Identical augmentations are applied to both context frames and the target to
maintain temporal consistency:
\begin{itemize}
  \item Horizontal flip: $p = 0.5$.
  \item Vertical flip: $p = 0.1$.
  \item Random rotation: $\pm 5°$.
  \item Brightness jitter: $\pm 0.1$.
  \item Contrast jitter: $\pm 0.1$.
\end{itemize}

\section{Training Configuration Summary}
\label{sec:config_summary}

The complete set of training hyperparameters is given in
Table~\ref{tab:hyperparams}.

\begin{table}[H]
\centering
\caption{Diffusion model training hyperparameters.}
\label{tab:hyperparams}
\begin{tabular}{ll}
\toprule
\textbf{Parameter} & \textbf{Value} \\
\midrule
Context window ($k$)                  & 5 frames \\
Image resolution                       & $256 \times 256$ px \\
Batch size                             & 32 \\
Training epochs                        & 1000 \\
Learning rate                          & $1 \times 10^{-4}$ \\
LR warm-up epochs                      & 10 \\
LR schedule                            & Cosine annealing \\
Minimum LR                             & $1 \times 10^{-7}$ \\
Gradient clip                          & 1.0 \\
Diffusion timesteps ($T$)             & 1000 \\
Noise schedule                         & Linear ($\beta_1{=}10^{-4}$, $\beta_T{=}2{\times}10^{-2}$) \\
DDIM sampling steps ($S$)             & 150 \\
DDIM $\eta$                           & 0 (deterministic) \\
Residual prediction                    & Yes \\
$v$-prediction                         & No ($\varepsilon$-prediction) \\
Time embedding dim ($d_\text{time}$)  & 8 \\
Context embedding dim ($D$)           & 96 \\
Context summary channels ($c_\ctx$)   & 4 \\
$\ell_1$-$x_0$ weight ($\lambda_{x_0}$) & 0.02 \\
Laplacian weight ($\lambda_\text{Lap}$) & 0.05 \\
Transition loss weight                 & 0.5 \\
Min-SNR $\gamma$                       & 5.0 \\
EMA decay                              & 0.9995 \\
Context perturb probability            & 0.50 (newest frame) \\
Context perturb noise std              & 0.07 \\
\bottomrule
\end{tabular}
\end{table}

% ==============================================================================
\chapter{Experiments and Results}
\label{chap:experiments}
% ==============================================================================

\section{Evaluation Setup}
\label{sec:eval_setup}

All experiments are conducted on the held-out \textbf{test split}, which
contains 69 embryos not seen during training or model selection.
Test pairs are constructed using the same windowed scheme described in
\Cref{ssec:pair_construction}.

Inference uses DDIM with 150 steps, $\eta = 0$ (deterministic), and a fixed
random seed.
All images are generated at $256 \times 256$ resolution (the training
resolution).

\subsection{Metrics}
\label{ssec:metrics}

The following metrics are computed:
\begin{itemize}
  \item \textbf{MSE}: mean squared error between generated and ground-truth
        frames (lower is better).
  \item \textbf{MAE}: mean absolute error (lower is better).
  \item \textbf{SSIM}: structural similarity index~\cite{wang2004ssim},
        which evaluates luminance, contrast, and structural similarity
        jointly (higher is better; perfect score = 1).
  \item \textbf{LPIPS}: learned perceptual image patch similarity~\cite{zhang2018lpips},
        based on deep feature distances (lower is better; perceptually
        perfect = 0).
\end{itemize}

\section{One-Step Next-Frame Prediction}
\label{sec:onestep}

The primary experiment evaluates one-step prediction under \emph{teacher
forcing}: the context window always contains ground-truth frames.
The model generates $\hat{x}_{t+1}$ for each test pair and the result is
compared against the ground-truth target.

Table~\ref{tab:onestep} reports the quantitative results averaged over all
test pairs.

\begin{table}[H]
\centering
\caption{Quantitative results for one-step next-frame prediction on the test
         set (teacher forcing).}
\label{tab:onestep}
\begin{tabular}{lcccc}
\toprule
\textbf{Setting} & \textbf{MSE} & \textbf{MAE} & \textbf{SSIM} & \textbf{LPIPS} \\
\midrule
One-step (all) & 0.1451 & 0.2472 & 0.4360 & 0.1165 \\
\bottomrule
\end{tabular}
\end{table}

Figure~\ref{fig:onestep_qual} shows representative one-step prediction
examples, with context frames, ground-truth target, and generated next frame.

\begin{figure}[H]
\centering
% \includegraphics[width=\textwidth]{figures/onestep_examples}
\caption{Qualitative one-step prediction examples. From left to right:
         context frames $x_{t-k+1}, \ldots, x_t$, ground truth $x_{t+1}$,
         and the model's prediction $\hat{x}_{t+1}$.}
\label{fig:onestep_qual}
\end{figure}

\section{Transition vs.\ Non-Transition Analysis}
\label{sec:transition}

To assess model behaviour at morphological transitions, test pairs are
partitioned into \emph{transition} pairs (phase label changes between $t$
and $t{+}1$) and \emph{non-transition} pairs.

Table~\ref{tab:transition} shows that performance is comparable in both
settings, suggesting that the model does not collapse when predicting across
phase boundaries.

\begin{table}[H]
\centering
\caption{One-step prediction performance on transition vs.\ non-transition
         test samples.}
\label{tab:transition}
\begin{tabular}{lcccc}
\toprule
\textbf{Subset} & \textbf{MSE} & \textbf{MAE} & \textbf{SSIM} & \textbf{LPIPS} \\
\midrule
Transition     & 0.1424 & 0.2431 & 0.4355 & 0.1181 \\
Non-transition & 0.1418 & 0.2445 & 0.4312 & 0.1100 \\
\bottomrule
\end{tabular}
\end{table}

\section{Multi-Step Rollout Evaluation}
\label{sec:rollout}

To evaluate temporal stability, the model is tested under \textbf{free-running
rollouts}: starting from a ground-truth context window, generated frames are
appended to the context and used as input for subsequent steps (no teacher
forcing).

Metrics are computed independently at each rollout depth by comparing
generated frames to the corresponding ground-truth frames.
Table~\ref{tab:rollout} and Figure~\ref{fig:rollout_metrics} present the
results.

\begin{table}[H]
\centering
\caption{Rollout performance as a function of prediction horizon (free-running,
         no teacher forcing). All four metrics degrade with rollout depth.}
\label{tab:rollout}
\begin{tabular}{lcccc}
\toprule
\textbf{Horizon} & \textbf{MSE} & \textbf{MAE} & \textbf{SSIM} & \textbf{LPIPS} \\
\midrule
$t{+}1$ & 0.1446 & 0.2376 & 0.4478 & 0.0866 \\
$t{+}2$ & 0.1891 & 0.2915 & 0.3572 & 0.1355 \\
$t{+}3$ & 0.1850 & 0.2921 & 0.3723 & 0.1678 \\
$t{+}4$ & 0.1931 & 0.3040 & 0.3464 & 0.1959 \\
$t{+}5$ & 0.2182 & 0.3307 & 0.3026 & 0.2131 \\
$t{+}6$ & 0.2491 & 0.3645 & 0.2501 & 0.2597 \\
$t{+}7$ & 0.2283 & 0.3491 & 0.2943 & 0.2569 \\
\bottomrule
\end{tabular}
\end{table}

\begin{figure}[H]
\centering
% \includegraphics[width=\textwidth]{figures/rollout_metrics}
\caption{Rollout SSIM and LPIPS as a function of prediction horizon.
         Performance degrades monotonically (with minor fluctuations),
         reflecting error accumulation in the autoregressive context.}
\label{fig:rollout_metrics}
\end{figure}

\begin{figure}[H]
\centering
% \includegraphics[width=\textwidth]{figures/rollout_examples}
\caption{Qualitative multi-step rollout example. Each row shows the
         ground-truth sequence (top) and the model's free-running predictions
         (bottom) for a single test embryo.}
\label{fig:rollout_qual}
\end{figure}

\section{Qualitative Observations}
\label{sec:qualitative}

Qualitative inspection of generated frames reveals that:
\begin{itemize}
  \item For short horizons ($t{+}1$, $t{+}2$), generated frames preserve the
        coarse embryo structure, cell boundary positions, and overall intensity
        distribution.
  \item At medium horizons ($t{+}3$--$t{+}5$), fine morphological detail
        starts to blur and some structural deformation becomes apparent.
  \item At longer horizons ($t{+}6$--$t{+}7$), noticeable structural drift
        and loss of cell boundary sharpness are observed.
  \item The model never collapses to a blurry mean; generated frames retain
        embryo-like appearance throughout the evaluated rollout window.
\end{itemize}

% ==============================================================================
\chapter{Discussion}
\label{chap:discussion}
% ==============================================================================

\section{Summary of Findings}
\label{sec:summary}

The experimental results demonstrate that the proposed context-conditioned
diffusion model can generate visually plausible next-frame predictions for
human embryo time-lapse sequences.
For short-term prediction (one step, teacher forcing), the model achieves
SSIM $= 0.436$ and LPIPS $= 0.116$, preserving coarse embryo structure and
maintaining temporal continuity with the conditioning context.
Transition and non-transition pairs show comparable performance, indicating
that the model does not systematically fail at morpho-kinetic boundaries.

Free-running rollout performance degrades with horizon, which is expected in
any autoregressive system subject to error accumulation.
The degradation is gradual rather than catastrophic, suggesting that the
context perturbation training scheme and the residual prediction
parameterisation are effective in bridging the train–rollout distribution
gap to some extent.

\section{Limitations}
\label{sec:limitations}

Several limitations affect the current approach:

\begin{enumerate}
  \item \textbf{Single-plane grayscale imagery.} The model is trained and
        evaluated on F$0$ only.
        The three-dimensional structure of the embryo is not captured, and
        focal-plane information that could inform depth-dependent morphological
        changes is ignored.

  \item \textbf{Pixel-level evaluation metrics.} MSE, MAE, SSIM, and LPIPS
        measure pixel and perceptual similarity but do not directly evaluate
        biological correctness.
        A generated frame with excellent SSIM may still correspond to a
        morphologically impossible developmental state.

  \item \textbf{No explicit phase supervision in the diffusion objective.}
        Phase annotations are available but are not incorporated into the
        training loss.
        The model does not explicitly reason about discrete developmental
        states, which may contribute to inconsistencies at phase transitions
        and over long rollouts.

  \item \textbf{Limited training scale.} Due to computational and time
        constraints within the thesis period, training duration, architectural
        search, and hyperparameter optimisation were limited.
        Longer training and larger models are expected to improve performance
        substantially.

  \item \textbf{Autoregressive error accumulation.} The model is trained
        with a one-step objective; there is no explicit long-horizon
        consistency constraint.
        This is the fundamental cause of rollout degradation and would require
        architectural or training-objective changes to address systematically.
\end{enumerate}

\section{Future Work}
\label{sec:future_work}

Building on the baseline established in this thesis, the following directions
are identified as the highest priority for future work:

\begin{enumerate}
  \item \textbf{Phase-conditioned or phase-aware training.}
        Incorporating phase labels into the diffusion training objective
        (e.g.\ via classifier-free guidance or phase-predictive auxiliary
        losses) may improve morphological consistency at transitions.

  \item \textbf{Multi-plane or volumetric conditioning.}
        Conditioning on multiple focal planes simultaneously would give the
        model access to 3D structural information and likely improve biological
        realism.

  \item \textbf{Long-horizon training objectives.}
        Curriculum strategies~\cite{bengio2009curriculum} or scheduled sampling
        of self-generated context during training could close the
        train/rollout distribution gap more aggressively than the current
        perturbation scheme.

  \item \textbf{Temporal consistency losses.}
        Auxiliary losses that penalise temporal inconsistency across multiple
        generated steps (e.g.\ via a temporal discriminator) could enforce
        smoother long-rollout trajectories.

  \item \textbf{Larger architectures.}
        The current model uses relatively modest channel widths
        $(32, 64, 128, 256, 256)$.
        Increasing model capacity with the same training recipe is likely to
        improve one-step and rollout performance.

  \item \textbf{Biologically-grounded evaluation.}
        Developing evaluation protocols that assess morpho-kinetic plausibility
        (e.g.\ whether generated sequences follow the correct phase order) would
        complement the current pixel-level metrics.
\end{enumerate}

% ==============================================================================
\chapter{Conclusion}
\label{chap:conclusion}
% ==============================================================================

This thesis presented a diffusion-based approach to next-frame prediction in
human embryo time-lapse sequences.
The key contributions are:
\begin{itemize}
  \item A reproducible preprocessing pipeline that produces a high-quality
        silver dataset from 704 embryos across 7 focal planes.
  \item A multi-class morpho-kinetic phase classifier (weighted F1$= 0.643$
        on F$0$) trained for analysis and optional rollout guidance.
  \item A context-conditioned DDPM model with cross-attention context
        conditioning, absolute frame-index temporal embedding, residual
        prediction parameterisation, and context perturbation for exposure bias
        mitigation.
  \item A systematic quantitative and qualitative evaluation covering
        one-step teacher-forced prediction, transition/non-transition analysis,
        and multi-step free-running rollouts.
\end{itemize}

The results demonstrate the feasibility of diffusion-based embryo time-lapse
prediction, with SSIM $= 0.436$ and LPIPS $= 0.116$ for one-step prediction
and plausible short-horizon rollouts.
Rollout quality degrades over longer horizons, motivating future work on
phase-aware conditioning, multi-horizon training objectives, and
biologically-grounded evaluation.

% ==============================================================================
%  Bibliography
% ==============================================================================
\printbibliography[heading=bibintoc, title={References}]

% ==============================================================================
%  Appendix
% ==============================================================================
\appendix

\chapter{Phase CSV Format}
\label{app:phase_csv}

Each embryo's annotation file follows the format:
\begin{verbatim}
phase,start_frame
pPB2,1
pPNa,4
pPNf,15
p2,22
...
\end{verbatim}
The \texttt{start\_frame} column records the first frame index at which
the corresponding phase was observed.
Phases without a recorded start frame are absent from the file.
Frames are assigned a phase label by finding the latest phase whose
\texttt{start\_frame} does not exceed the frame index in question.

\chapter{Hough Circle Detection Parameters}
\label{app:hough}

Embryo-centred cropping uses OpenCV's \texttt{HoughCircles} function
applied to a Gaussian-blurred Canny edge map.
Parameters are tuned to the expected embryo size range in the dataset
(typically $200$--$350$ pixel diameter in the original $500 \times 500$
resolution).
When no circle is detected (fewer than 1 hit), the fallback centre crop
extracts a fixed $400 \times 400$ window from the image centre before
resizing.

\chapter{Repository Structure}
\label{app:repo}

\begin{verbatim}
thesis/
├── data/
│   ├── embryo_dataset_bronze/
│   ├── embryo_dataset_annotations/
│   └── embryo_dataset_silver/<plane>/<embryo>/
├── preprocessing/
│   ├── preprocess_dataset.ipynb
│   └── train_presence_classifier.ipynb
├── classifier/
│   └── train_phase_classifier.ipynb
├── diffusion/
│   ├── models.py
│   ├── train_diffusion.ipynb
│   ├── embryo_diffusion_models/
│   ├── diffusion_samples/
│   └── cache/
├── figures/
├── requirements.txt
└── README.md
\end{verbatim}

\end{document}
